\chapter{Conclusion}\label{chap:conclusion}

The conclusion summarizes the thesis and discusses future work.
Do not introduce any new material here.
A reader who reads only the introduction and the conclusion should get a coherent
picture of what was done and what remains open.
It should follow the same structure as the introduction and the abstract, 
but with more emphasis on the bottom part of the hourglass.

%-----------------------------------------------------------------------
% Summary
%-----------------------------------------------------------------------

The summary should address the following points:
\begin{enumerate}
    \item Restate the problem studied and explain why it is relevant.
    \item Describe the main contribution and the approach taken
          (reference \cref{chap:main-chapter}).
    \item Name the central theoretical result (reference \cref{thm:main})
          and state informally what it guarantees.
    \item Summarize the experimental findings -- number of additional examples
          solved, runtime overhead -- if an evaluation was conducted
          (reference \cref{chap:evaluation}).
\end{enumerate}

%-----------------------------------------------------------------------
% Future Work
%-----------------------------------------------------------------------

Each future work item should:
\begin{enumerate}
    \item Identify a specific limitation of the current approach or a natural
          extension that was beyond the scope of this thesis.
    \item Describe what would be needed to address it (new theoretical machinery,
          a different encoding, a richer formalism, etc.).
    \item Explain what benefit this would bring (more examples solved, a stronger
          guarantee, applicability to a wider class of systems).
\end{enumerate}